\smallskip
\emph{Linear and unrestricted types.}
Our 
%formal 
calculus is mostly linear, except that we
%. The only exception we make is to
allow unrestricted bindings for recursive functions, which would be
rendered useless otherwise. 
%On the other hand, we want to speak about
%recursive protocols, so we felt like recursive functions are needed in
%the formal development. 
Our implementation goes well beyond that in
allowing fully fledged unrestricted computations besides ensuring that
session typed channels are treated linearly.
% That said, all related work is closer to our implementation than to
% the formal calculus.
%
There is some work on recursion and linearity in the context of lambda calculus
\cite{DBLP:journals/lisp/AlvesFFM10}
%. This work 
adding iterators on
linear natural numbers to the calculus, thereby avoiding the need for
a rule like {\rulenameexpUVar}.
%, which we need to bind the recursive function to a name.

System F\degree~\cite{DBLP:conf/tldi/MazurakZZ10}
extends System F with kinds to integrate polymorphism with linear and
unrestricted types. This work serves as a blueprint for much
subsequent work in this area including LFST and FreeST (see below), as well as our
implemented language.
They categorize all objects, including functions, into one-use or
many-use resources. This style can be traced back to 
\citet{walker:substructural-type-systems}. 
% One example in the paper applies the language to protocols built with
% linear capabilities, reminiscent of session types.
%
LFST \cite{lindley17:_light_funct_session_types} formalizes the
session type system of Links, which features a mix of linear and
unrestricted types, embedding a standard functional programming
language. The substructural behavior of types is specified by a
sophisticated kind structure.  The kinding used in our implementation
is inspired by LFST, though LFST has additional kinds to describe row
types that we do not support.
%
\citet{DBLP:journals/corr/abs-2106-06658}
consider an integration of context-free session types with
System~F. Their type system includes a kinding discipline that manages
linear and unrestricted versions of ordinary types and
message types (possible payloads for the send and receive
operations). Our implemented system extends the kind structure in a
similar way.

Linear Haskell (LH) \cite{DBLP:journals/pacmpl/BernardyBNJS18}
is a proposal to integrate linear types with stock functional
programming. While the work discussed so far characterizes resources,
LH supports the lollipop type $S \multimap T$ which classifies
functions that use their argument exactly once. Originally, LH forced
the programmer to allocate resources using continuations. This restriction
has been fixed in subsequent work \cite{DBLP:journals/pacmpl/SpiwackKBWE22}.
%
There is further work on integrating linear and dependent types (e.g.,
\cite{DBLP:conf/popl/KrishnaswamiPB15}) which relies on
bifurcating the language in two calculi connected by an adjunction.

% We believe that \algst definitions that do not make
% significant use of negation on protocols

%%% Local Variables:
%%% mode: latex
%%% TeX-master: "main"
%%% End:
